\documentclass[12pt]{article}

\usepackage{fullpage}
\usepackage{amssymb}
\usepackage{amsmath}

\newtheorem{theorem}{Theorem}
\newtheorem{lemma}[theorem]{Lemma}
\newtheorem{proposition}[theorem]{Proposition}
\newtheorem{corollary}[theorem]{Corollary}
\theoremstyle{definition}
\newtheorem{definition}[theorem]{Definition}
\newtheorem{remark}[theorem]{Remark}

\begin{document}

\title{The conjectured order-10 recurrence for the number of
nonequivalent maximum king placements on a $4 \times 2n$ board}

\author{Shi Li\thanks{Business School, Nanjing Normal University, Nanjing,
China. Email: \texttt{saberson946@gmail.com}.}}

\date{}

\maketitle

\begin{abstract}
Let $a(n)$ denote the number of ways to place $2n$ nonattacking kings on a
$4 \times 2n$ board attaining the maximum $2n$, counted up to the symmetry
group of the rectangle. In 2018 Colin Barker conjectured an order-10 linear
recurrence for $a(n)$, equivalently a rational generating function, on the
basis of about eleven confirming terms; the conjecture is recorded as such in
the OEIS entry A321614 and has remained unproved. We prove it for all $n$. The
proof has two ingredients. First, the maximum condition forces exactly one king
into each of the $2n$ disjoint $2 \times 2$ blocks of the board, which turns the
problem into a walk in a $12$-state transfer graph. Second, Burnside's lemma
over the Klein four-group expresses $a(n)$ as an average of four fixed-point
counts, each of which is a linear combination of entries of powers of the
transfer matrix. Summing the resulting rational generating functions in exact
arithmetic yields Barker's generating function identically. Every combinatorial
step is proved by hand; the only machine-assisted step is the exact rational
arithmetic of the final identity, which the accompanying programs reproduce in
seconds.
\end{abstract}

\section{Introduction}

A set of kings on a chessboard is \emph{nonattacking} if no two of them are at
Chebyshev distance at most one, that is, if no two share a row, column, or
diagonal neighbour. On a $4 \times 2n$ board the maximum number of nonattacking
kings is $2n$, since the board partitions into $2n$ disjoint $2 \times 2$ blocks
and each block accommodates at most one king.

Let $a(n)$ be the number of such maximum placements, counted up to the
symmetry group of the rectangle, that is, up to the Klein four-group generated
by the horizontal reflection, the vertical reflection, and the $180^\circ$
rotation. This is the sequence A321614 in the OEIS \cite{oeis}. Its published
terms are
\[
1, 4, 23, 106, 473, 1939, 7618, 28703, 105112, 375597, 1316944, 4544124,
\]
continuing to $a(21) = 665017743665$.

On 22 December 2018 Colin Barker conjectured \cite{oeis} that
\begin{equation}
\label{eq:rec}
a(n) = 12a(n-1) - 54a(n-2) + 98a(n-3) + 17a(n-4) - 346a(n-5)
 + 505a(n-6) - 210a(n-7) - 120a(n-8) + 126a(n-9) - 27a(n-10)
\end{equation}
for $n > 9$, with the equivalent rational generating function
\begin{equation}
\label{eq:gf}
A(x) = \sum_{n \ge 0} a(n) x^n
 = \frac{(1-2x)(1-6x+17x^2-18x^3-2x^4+7x^5+6x^6-3x^7)}
        {(1-x)^2 (1-3x)^2 (1-3x+x^2)(1-x-x^2)(1-3x^2)} .
\end{equation}
The evidence was a fit to roughly eleven terms. We prove the conjecture.

\begin{theorem}
\label{thm:main}
For every $n \ge 10$ the recurrence \eqref{eq:rec} holds, and the generating
function of $a(n)$ is \eqref{eq:gf}.
\end{theorem}

\begin{remark}
\label{rem:group}
For $n = 2$ the board $4 \times 4$ is a square and its full symmetry group is
larger than the Klein four-group. The sequence is nevertheless defined with the
four-group for every $n$, and it is this convention that reproduces
$a(2) = 23$. The distinction matters only at $n = 2$ and is flagged because it
is the one place where the definition and the geometry diverge.
\end{remark}

The proof strategy is standard for counts of this type: a finite-state transfer
matrix for the placements, Burnside's lemma for the symmetry quotient, and exact
rational arithmetic for the generating function. The point of the paper is that
the maximum condition makes the transfer matrix small and explicit, so that the
whole computation can be carried out and checked.

\section{The block structure}

Index board rows by $0, 1, 2, 3$ and columns by $0, 1, \dots, 2n-1$. Partition
the board into blocks
\[
B_{i,j} = \{2i, 2i+1\} \times \{2j, 2j+1\},
\qquad i \in \{0,1\}, \quad j \in \{0, \dots, n-1\}.
\]
There are $2n$ blocks.

\begin{lemma}
\label{lem:pigeonhole}
Every placement of $2n$ nonattacking kings puts exactly one king in each block
$B_{i,j}$.
\end{lemma}

\begin{proof}
Two cells of one block differ by at most one in each coordinate, so they are at
Chebyshev distance at most one and a block contains at most one king. Since
there are $2n$ blocks and $2n$ kings, every block contains exactly one.
\end{proof}

Consequently every maximum placement is described by a word
$s_0 s_1 \cdots s_{n-1}$ of \emph{column states}: if the king in block $B_{i,j}$
occupies the cell $(2i + r_{i,j},\, 2j + c_{i,j})$ with $r_{i,j}, c_{i,j} \in
\{0,1\}$, then
\[
s_j = (r_{0,j},\, c_{0,j},\, r_{1,j},\, c_{1,j}).
\]

\begin{lemma}
\label{lem:locality}
Two blocks can interact only if they lie in the same column or in adjacent
columns.
\end{lemma}

\begin{proof}
Block $B_{i,j}$ occupies board columns $2j$ and $2j+1$. If $|j - j'| \ge 2$
then any two cells of $B_{i,j}$ and $B_{i',j'}$ have column indices differing
by at least $4 - 1 = 3$, so the kings do not attack.
\end{proof}

By Lemma \ref{lem:locality} every constraint relates a column to itself or to
one of its two neighbours, and no constraint relates cells within a block,
by Lemma \ref{lem:pigeonhole}.

\section{The four local constraints}

Write $s = (r_0, c_0, r_1, c_1)$ for a column state and $u = (R_0, C_0, R_1,
C_1)$ for the state of the next column.

\subsection{Within a column}

Blocks $B_{0,j}$ and $B_{1,j}$ occupy the same board columns, so their column
indices differ by at most one and the column condition is automatic. The
vertical distance between the kings is
\[
(2 + R_1) - r_0 = 2 + R_1 - r_0,
\]
which is at most one exactly when $r_0 = 1$ and $R_1 = 0$. Hence $s$ is
\emph{admissible}, meaning that the two kings of the column do not attack, if
and only if
\begin{equation}
\label{eq:adm}
\neg (r_0 = 1 \ \text{and}\ r_1 = 0).
\end{equation}
Among the sixteen states in $\{0,1\}^4$ exactly twelve satisfy \eqref{eq:adm}.

\subsection{Same row-pair, adjacent columns}

Blocks $B_{i,j}$ and $B_{i,j+1}$ occupy the same pair of rows. The horizontal
distance between the kings is
\[
(2j + 2 + C) - (2j + c) = 2 + C - c,
\]
which is at most one exactly when $c = 1$ and $C = 0$. Hence the condition is
\begin{equation}
\label{eq:same}
c_{i,j} \le c_{i,j+1}
\end{equation}
for $i = 0$ and for $i = 1$.

\subsection{Diagonal, first orientation}

Consider $B_{0,j}$ and $B_{1,j+1}$. The vertical distance is $2 + R_1 - r_0$,
at most one exactly when $r_0 = 1$ and $R_1 = 0$. The horizontal distance is
$2 + C_1 - c_0$, at most one exactly when $c_0 = 1$ and $C_1 = 0$. The kings
attack precisely when all four hold, so the condition is
\begin{equation}
\label{eq:diag1}
\neg (r_0 = 1 \ \text{and}\ R_1 = 0 \ \text{and}\ c_0 = 1 \ \text{and}\ C_1 = 0).
\end{equation}

\subsection{Diagonal, second orientation}

Symmetrically, for $B_{1,j}$ and $B_{0,j+1}$ the condition is
\begin{equation}
\label{eq:diag2}
\neg (r_1 = 0 \ \text{and}\ R_0 = 1 \ \text{and}\ c_1 = 1 \ \text{and}\ C_0 = 0).
\end{equation}

By Lemmas \ref{lem:pigeonhole} and \ref{lem:locality}, conditions
\eqref{eq:adm} to \eqref{eq:diag2} are necessary and sufficient for a word of
column states to describe a maximum nonattacking placement. This is a finite
case analysis, complete by the four derivations above.

\section{The transfer matrix}

List the twelve admissible states in the order
\[
\begin{array}{llll}
s_0 = (0,0,0,0) & s_1 = (0,0,0,1) & s_2 = (0,0,1,0) & s_3 = (0,0,1,1) \\
s_4 = (0,1,0,0) & s_5 = (0,1,0,1) & s_6 = (0,1,1,0) & s_7 = (0,1,1,1) \\
s_8 = (1,0,1,0) & s_9 = (1,0,1,1) & s_{10} = (1,1,1,0) & s_{11} = (1,1,1,1)
\end{array}
\]
and let $M$ be the $12 \times 12$ matrix with $M[i][j] = 1$ if the pair
$(s_i, s_j)$ satisfies \eqref{eq:same}, \eqref{eq:diag1} and \eqref{eq:diag2},
and $M[i][j] = 0$ otherwise. Explicitly,
\[
M = \begin{pmatrix}
1&1&1&1&1&1&1&1&1&1&1&1\\
0&1&0&1&0&1&0&1&0&0&0&1\\
1&1&1&1&1&1&1&1&1&1&1&1\\
0&1&0&1&0&1&0&1&0&1&0&1\\
0&0&0&0&1&1&1&1&0&0&1&1\\
0&0&0&0&0&1&0&1&0&0&0&1\\
0&0&0&0&1&1&1&1&0&0&1&1\\
0&0&0&0&0&1&0&1&0&0&0&1\\
1&1&1&1&1&1&1&1&1&1&1&1\\
0&1&0&1&0&1&0&1&0&1&0&1\\
0&0&0&0&0&1&1&1&0&0&1&1\\
0&0&0&0&0&1&0&1&0&0&0&1
\end{pmatrix}.
\]
Let $\mathbf{1}$ denote the all-ones vector. A word of $n$ admissible states
with every consecutive pair satisfying the transition conditions is exactly a
maximum placement, so the number of placements without any symmetry quotient is
\begin{equation}
\label{eq:T}
T(n) = \mathbf{1}^{\mathsf{T}} M^{n-1} \mathbf{1},
\qquad T(0) = 1.
\end{equation}
The characteristic polynomial of $M$ is
\begin{equation}
\label{eq:charpoly}
\chi_M(\lambda) = \lambda^6 (\lambda - 3)^2 (\lambda^2 - 3\lambda + 1)^2,
\qquad
\det(I - xM) = (3x-1)^2 (x^2 - 3x + 1)^2 .
\end{equation}

\section{Burnside over the four-group}

The four symmetries act on cells by
\[
h \colon (\rho, \gamma) \mapsto (\rho, 2n-1-\gamma), \quad
v \colon (\rho, \gamma) \mapsto (3-\rho, \gamma), \quad
r \colon (\rho, \gamma) \mapsto (3-\rho, 2n-1-\gamma).
\]
On blocks they act by $h \colon B_{i,j} \mapsto B_{i, n-1-j}$,
$v \colon B_{i,j} \mapsto B_{1-i, j}$ and $r \colon B_{i,j} \mapsto
B_{1-i, n-1-j}$, and on offsets by $h \colon (r,c) \mapsto (r, 1-c)$,
$v \colon (r,c) \mapsto (1-r, c)$ and $r \colon (r,c) \mapsto (1-r, 1-c)$. The
induced maps on column states are therefore
\begin{align}
\sigma_h (r_0,c_0,r_1,c_1) &= (r_0, 1-c_0, r_1, 1-c_1), \label{eq:sigh}\\
\sigma_v (r_0,c_0,r_1,c_1) &= (r_1, c_1, r_0, c_0), \label{eq:sigv}\\
\sigma_r (r_0,c_0,r_1,c_1) &= (1-r_1, 1-c_1, 1-r_0, 1-c_0). \label{eq:sigr}
\end{align}
With the convention of Remark \ref{rem:group}, Burnside's lemma gives, for
$n \ge 1$,
\begin{equation}
\label{eq:burnside}
a(n) = \frac{T(n) + F_h(n) + F_v(n) + F_r(n)}{4},
\end{equation}
where $T(n)$ is as in \eqref{eq:T} and $F_h, F_v, F_r$ count the placements
fixed by $h$, $v$ and $r$ respectively.

\subsection{The vertical reflection}

By \eqref{eq:sigv}, a column state is fixed by $\sigma_v$ if and only if
\[
r_1 = 1 - r_0 \quad \text{and} \quad c_1 = c_0 .
\]
The admissibility condition \eqref{eq:adm} then reads $\neg(r_0 = 1 \ \text{and}
\ r_0 = 1)$, forcing $r_0 = 0$ and $r_1 = 1$. The fixed admissible states are
therefore exactly
\[
(0,c,1,c), \qquad c \in \{0,1\},
\]
namely $s_2$ and $s_7$. The two other states satisfying the algebraic fixed
condition, $(1,0,0,0)$ and $(1,1,0,1)$, are inadmissible and do not occur.

For two consecutive such states $(0,c,1,c)$ and $(0,C,1,C)$, condition
\eqref{eq:same} reads $c \le C$; conditions \eqref{eq:diag1} and \eqref{eq:diag2}
are vacuously satisfied because they require $r_0 = 1$ and $r_1 = 0$
respectively. Hence a placement fixed by $v$ is exactly a non-decreasing word in
$\{0,1\}$ of length $n$, and
\begin{equation}
\label{eq:Fv}
F_v(n) = n + 1 .
\end{equation}

\subsection{The horizontal reflection}

A placement fixed by $h$ has column $2h-1-j$ in state $\sigma_h(s_j)$, so for
even $n = 2h$ it is a palindrome determined by $s_0, \dots, s_{h-1}$. Every
transition inside a half is automatic once the word is admissible; the only
additional condition is at the junction between columns $h-1$ and $h$, which
reads
\[
\text{trans}\bigl(s_{h-1},\, \sigma_h(s_{h-1})\bigr).
\]
Let $J_h$ be the indicator vector of the states satisfying this condition. By
\eqref{eq:sigh} and the transition conditions,
\[
J_h = (1,0,1,0,\ 0,0,0,0,\ 1,0,0,0).
\]
Then $(M^{h-1} J_h)[s_0]$ counts the admissible words of length $h$ beginning at
$s_0$ whose last letter satisfies the junction condition, so
\begin{equation}
\label{eq:Fh}
F_h(2h) = \mathbf{1}^{\mathsf{T}} M^{h-1} J_h,
\qquad h \ge 1 .
\end{equation}
For odd $n$ the middle column would have to satisfy $\sigma_h(s) = s$, which by
\eqref{eq:sigh} requires $c_0 = 1 - c_0$, impossible. Hence
\begin{equation}
\label{eq:Fhodd}
F_h(n) = 0 \quad \text{for odd } n .
\end{equation}

\subsection{The $180^\circ$ rotation}

The same analysis applies with \eqref{eq:sigr}. A state is fixed by $\sigma_r$
precisely when $r_0 = 1 - r_1$ and $c_0 = 1 - c_1$. Admissibility then forces
$r_0 = 0$ and $r_1 = 1$, leaving the two states
\[
s_3 = (0,0,1,1) \quad \text{and} \quad s_6 = (0,1,1,0).
\]
Write $C_1 = 2$ for their number. For even $n = 2h$ the junction condition gives
the indicator vector
\[
J_{re} = (1,0,1,1,\ 1,0,1,0,\ 1,1,0,0),
\]
and hence
\begin{equation}
\label{eq:Ffre}
F_r(2h) = \mathbf{1}^{\mathsf{T}} M^{h-1} J_{re}, \qquad h \ge 1 .
\end{equation}
For odd $n = 2h+1$ the middle column $m$ must be fixed by $\sigma_r$ and must
connect to its neighbour on both sides, so define
\[
J_{ro}[s] = \# \{\, m : \sigma_r(m) = m,\ \text{trans}(s,m),\ \text{trans}(m,\sigma_r(s)) \,\}.
\]
Then $F_r(1) = C_1 = 2$, and
\begin{equation}
\label{eq:Ffro}
F_r(2h+1) = \mathbf{1}^{\mathsf{T}} M^{h-1} J_{ro}, \qquad h \ge 1 .
\end{equation}

\section{The generating function}

Each of $T$, $F_h$ and $F_r$ is a sum of sequences of the form
$\mathbf{1}^{\mathsf{T}} M^{k} w$ over $k$, hence rational. By
Cayley--Hamilton and \eqref{eq:charpoly}, for any vectors $u$ and $w$,
\[
\sum_{k \ge 0} (u^{\mathsf{T}} M^k w) x^k = \frac{p(x)}{q(x)},
\qquad q(x) = \det(I - xM),
\]
with $\deg p < \deg q$. The numerator is recovered as the initial segment of
$q(x)$ times the series, an identity valid because the coefficients satisfy the
recurrence implied by $q$ beyond degree $\deg q$. The vanishing of $F_h$ at odd
$n$ in \eqref{eq:Fhodd} and the even and odd cases in \eqref{eq:Fh},
\eqref{eq:Ffre} and \eqref{eq:Ffro} correspond to the substitution $y = x^2$,
and \eqref{eq:Fv} contributes
\[
\sum_{n \ge 1} (n+1) x^n = \frac{1}{(1-x)^2} - 1 .
\]
All of this is carried out in exact rational arithmetic. Summing the four pieces
through \eqref{eq:burnside} and cancelling gives
\[
\sum_{n \ge 0} a(n) x^n
 = \frac{6x^8 - 15x^7 - 8x^6 + 11x^5 + 34x^4 - 52x^3 + 29x^2 - 8x + 1}
        {27x^{10} - 126x^9 + 120x^8 + 210x^7 - 505x^6 + 346x^5
         - 17x^4 - 98x^3 + 54x^2 - 12x + 1},
\]
where the denominator factors as
\[
(x-1)^2 (3x-1)^2 (3x^2-1) (x^2-3x+1) (x^2+x-1),
\]
which is exactly the denominator of \eqref{eq:gf}, and the numerator is exactly
that of \eqref{eq:gf}. As the two generating functions are equal as rational
functions, and a rational function with denominator of degree $10$ satisfies the
corresponding order-$10$ recurrence, \eqref{eq:rec} holds for all $n \ge 10$.
This proves Theorem \ref{thm:main}. The verification is reproduced by the
accompanying programs, which compute the difference of the two rational
functions symbolically and report zero.

\section{Verification}

The programs accompanying this paper are collected in a single directory and
are runnable with a standard Python installation. They perform the following
checks.

\begin{enumerate}
\item The transfer matrix reproduces all 22 published terms of A321614 exactly.

\item $T(n)$, $F_h(n)$, $F_v(n)$ and $F_r(n)$ are each checked \emph{separately}
against definition-level enumeration of placements and of fixed points, for
$n \le 5$. Checking the four terms individually rather than only their average
matters: Burnside's lemma is a statement about each fixed-point count on its
own, so two terms could be wrong in opposite directions and cancel in
\eqref{eq:burnside}.

\item An independent implementation that never uses the block structure,
namely a column-profile dynamic program over grid independent sets, agrees with
the transfer matrix for $n \le 26$.

\item The generating-function identity is verified symbolically, and the
reported difference of the two rational functions is exactly zero.

\item The objects transcribed in this paper, namely the matrix $M$, the state
list, and the vectors $J_h$ and $J_{re}$, are parsed back out of the source of
this paper and compared with the programs, so that a transcription error cannot
go unnoticed.

\item The dimension reduction behind Lemma \ref{lem:pigeonhole} is checked
exhaustively: all 16 within-column states and all 144 ordered pairs of
cross-column states are compared with the geometric nonattacking condition.

\end{enumerate}

The programs, the term data, and the instructions for reproducing every number
in this paper are available from the author.

\section{Acknowledgment}

The author thanks the OEIS for maintaining the entry A321614 and Colin Barker
for the conjecture, and acknowledges the collaborative superpermutation and
nonattacking-pieces literature from which the transfer-matrix method used here
is drawn.

\begin{thebibliography}{9}

\bibitem{oeis}
N. J. A. Sloane, editor, \emph{The On-Line Encyclopedia of Integer Sequences},
sequence A321614, \url{https://oeis.org/A321614}.

\bibitem{howroyd}
A. Howroyd, comment on sequence A321614, OEIS, December 2018.

\bibitem{kotesovec}
V. Kotesovec, \emph{Non-attacking chess pieces}, sixth edition, 2013.

\bibitem{wilf}
H. S. Wilf, The problem of the kings, \emph{Electronic Journal of
Combinatorics} \textbf{2} (1995), R3.

\end{thebibliography}

\end{document}
